\section{Reducedness of the actual hypersurface transform}
\label{sec:hironaka-hypersurface-strict-ideal}

The maximum-multiplicity centre construction produces exact exceptional
multiplicity.  The remaining hypersurface issue concerns the entire residual
ideal, including its reducedness at points of the exceptional divisor.
The following local analytic argument detects a repeated factor on the same
genuine regular piece as the tests excluding the exceptional divisor.
The general current-stage centre assertion of
\cite{bierstone-milman97-invent-math}, Theorem~1.6(1), is a separate objective.

\begin{lemma}[Reducedness across a divisor without a common factor]
\label{hir:reduced-germ-across-divisor}
\lean{Hironaka.HironakaHypersurfaceStrictIdeal.squarefree_germ_of_reduced_off_divisor}
\leanok
Let $f,e$ be holomorphic germs at the origin of $\mathbb C^{m+1}$, with
$f\ne0$.  Suppose no prime factor of $f$ divides $e$.  On a neighbourhood
of the origin, suppose that at every $x$ with $e(x)\ne0$, the actual
principal ideal germ $(f)_x$ is radical.  Then the original germ $f$ is
squarefree.
\end{lemma}
\begin{proof}
The actual holomorphic germ ring is factorial.  Suppose a prime $p$ occurs
at least twice in $f$.  Its finite maximum exponent gives
$f=p^k b$, where $k\ge2$ and $p\nmid b$.  Choose holomorphic
representatives of these original germs; the identity holds on a
neighbourhood of the origin.  Generic regular differential avoidance for
the original prime, with the simultaneous tests $e$ and $b$, gives points
arbitrarily near the origin at which $p(x)=0$, $e(x)b(x)\ne0$, and the
differential of the same representative $p$ is nonzero in a fixed direction.
Choose such a point inside the neighbourhood where the factorization and
the off-divisor radicality assertions hold.

Since $b(x)\ne0$, its reciprocal is holomorphic near $x$.  Thus
$p^k=b^{-1}f$ belongs to $(f)_x$.  Radicality implies $p\in(f)_x$,
so $p=af$ for a holomorphic germ $a$ at $x$.  On the other hand,
$f=p^k b$ and $k\ge2$ imply both $f(x)=0$ and $df_x=0$.
Differentiating $p=af$ gives $dp_x=0$, contrary to the selected
nonzero differential.  Hence every prime has exponent at most one and
$f$ is squarefree.  If $g^r\in(f)$ for $r>0$, squarefree divisibility
implies $f\mid g$, which proves radicality of the principal ideal.
\end{proof}

\begin{lemma}[A nonexceptional equation cannot be supported only on the divisor]
\label{hir:nonexceptional-equation-not-divisor-supported}
\lean{Hironaka.HironakaHypersurfaceStrictIdeal.ne_zero_at_zero_of_zero_locus_subset_coordinate}
\leanok
Let $f$ be a nonzero actual holomorphic germ at the origin, and suppose
the coordinate germ $z_j$ does not divide it.  If the zero set of $f$
is contained in $\{z_j=0\}$ near the origin, then $f(0)\ne0$.
\end{lemma}
\begin{proof}
If $f(0)=0$, its germ is a nonunit.  Actual factoriality supplies a prime
factor $p$ of $f$.  The coordinate prime cannot be associated to $p$,
since otherwise it would divide $f$.  Therefore $p\nmid z_j$.
Genuine prime avoidance produces zeros of $p$ arbitrarily near the origin
with $z_j\ne0$.  Choose such a point inside a neighbourhood on which
the actual factorization $f=pb$ and the given zero-set containment both
hold.  There $f=0$ while $z_j\ne0$, a contradiction.
\end{proof}

\begin{lemma}[A whole-ideal witness excludes the coordinate factor]
\label{hir:ideal-nonexceptional-witness-coordinate-factor}
\lean{Hironaka.HironakaHypersurfaceStrictIdeal.not_coordinate_dvd_of_nonexceptional_witness}
\leanok
Let an actual ideal stalk $R_0$ be generated by a holomorphic germ $f$.
If $R_0$ has an actual member outside the vanishing ideal of the coordinate
hyperplane $\{z_j=0\}$, then $z_j$ does not divide $f$.
\end{lemma}
\begin{proof}
If $f=z_j b$ as holomorphic germs, then $f$ vanishes on that hyperplane
near the origin.  Every holomorphic multiple of $f$, and every finite
holomorphic linear combination of such multiples, also vanishes there.
Thus the entire original stalk $R_0$ lies in the actual hyperplane
vanishing ideal.  This contradicts the supplied actual member.
\end{proof}

\begin{corollary}[The actual coordinate-divisor criterion]
\label{hir:reduced-germ-across-coordinate-divisor}
\uses{hir:reduced-germ-across-divisor}
\lean{Hironaka.HironakaHypersurfaceStrictIdeal.squarefree_germ_of_reduced_off_coordinate}
\leanok
In the preceding lemma, let $e(z)=z_j$.  It suffices to assume that the
actual coordinate germ $z_j$ does not divide $f$.  Then $f$ is squarefree.
\end{corollary}
\begin{proof}
The coordinate germ is prime: its nonzero coordinate derivative rules out
a factorization into two nonunits, and the actual factorial germ ring
turns irreducibility into primeness.  A prime factor of $f$ dividing
$z_j$ would therefore be associated to $z_j$.  This would make $z_j$
divide $f$, contrary to the hypothesis.  Apply the preceding lemma.
\end{proof}

\begin{lemma}[The full ideal of a squarefree equation]
\label{hir:squarefree-equation-full-vanishing-ideal}
\lean{Hironaka.HironakaHypersurfaceStrictIdeal.vanishingIdeal_eq_span_of_squarefree}
\leanok
For a nonzero squarefree actual holomorphic germ $f$ at the origin,
the entire vanishing ideal of its zero germ is the principal ideal $(f)$.
\end{lemma}
\begin{proof}
The actual radical Nullstellensatz and factoriality give a nonzero
squarefree germ $r$ with the same zero germ and with full vanishing ideal
$(r)$.  The reduced-equation unit-transition theorem applied to the two
actual squarefree equations gives $f=ur$ with $u$ a holomorphic unit.
Thus $f$ and $r$ divide each other, their principal ideals are equal,
and the whole vanishing ideal is $(f)$.
\end{proof}

\begin{lemma}[Invertibility of the actual controlled hypersurface ideal]
\label{hir:invertible-controlled-hypersurface-ideal}
\lean{Hironaka.HironakaHypersurfaceStrictIdeal.isInvertible_controlledTransform}
\leanok
Let $\pi:Y\to X$ be an actual blowup along a closed smooth centre $C$,
and let $I$ be an invertible holomorphic ideal with
$I\subseteq I_C^\mu$.  The actual controlled transform
$R=(\pi^*I:I_{\pi^{-1}C}^{\mu})$ is invertible.
\end{lemma}
\begin{proof}
The exceptional locus of an actual blowup has empty interior.  The pullback
of the invertible ideal is therefore an actual invertible ideal: its local
nonzero generator cannot vanish identically after this pullback, since
the blowup is locally biholomorphic off the exceptional divisor.
The exceptional divisor ideal and each of its powers are invertible.
The given centre divisibility, pulled back through the actual blowup,
places $\pi^*I$ inside the exceptional power.  The actual colon of nested
invertible ideals is invertible: divide their local generators, then use
density of the nonzero denominator locus to cancel it as germs.
The separate coherent colon-product identity supplies the usual
factorization when needed below.
\end{proof}

\begin{lemma}[The actual off-centre controlled-ideal dictionary]
\label{hir:controlled-hypersurface-off-centre-stalk}
\lean{Hironaka.HironakaHypersurfaceStrictIdeal.controlledTransform_vanishingIdeal_mem_iff_of_notMem}
\leanok
For any actual blowup $\pi:Y\to X$ along $C$, any closed set $W$, any
$\mu\ge0$, and any $y$ with $\pi(y)\notin C$, the entire stalk of
the actual controlled transform of $(I_W,\mu)$ is the actual reduced
vanishing ideal of $\pi^{-1}W$ at $y$.
\end{lemma}
\begin{proof}
The exceptional divisor ideal contains $1$ at $y$: the original centre
ideal contains $1$ near $\pi(y)$, and actual pullback preserves this
membership.  Induction on $\mu$ gives $1$ in every exceptional power.
Thus a member $g$ of the actual colon belongs to the original pullback
ideal by testing the colon condition against $1$.
The genuine holomorphic inverse of the blowup off its centre identifies
the complete pullback ideal stalk with $I_{\pi^{-1}W,y}$.
Conversely, if $g$ belongs to this vanishing ideal, multiplying it by
any holomorphic member of the exceptional power preserves its vanishing
on $\pi^{-1}W$.  The same genuine inverse puts each such product in
the original pullback ideal, which is exactly the actual colon condition.
This proves both directions for every holomorphic germ, not just for a
chosen generator or zero locus.
\end{proof}

\begin{lemma}[Support equality across a smooth divisor]
\label{hir:invertible-ideal-support-across-smooth-divisor}
\uses{hir:ideal-nonexceptional-witness-coordinate-factor,hir:nonexceptional-equation-not-divisor-supported}
\lean{Hironaka.HironakaHypersurfaceStrictIdeal.zeroLocus_eq_of_off_smooth_divisor}
\leanok
Let $R$ be an actual invertible ideal on a complex manifold of dimension
$m+1$, let $E$ be a closed smooth hypersurface, and let $S$ be closed
with $S\subseteq\overline{S\setminus E}$.  Suppose the whole ideal
stalks $R_x$ and $I_{S,x}$ agree for $x\notin E$, and that at each
$x\in E$ some actual member of $R_x$ lies outside $I_{E,x}$.
Then the actual zero locus of $R$ is exactly $S$.
\end{lemma}
\begin{proof}
The coherent zero locus of $R$ is closed, and contains $S\setminus E$
by the off-divisor whole-stalk equality.  Density therefore puts all of
$S$ in this zero locus.
Conversely, away from $E$ the off-divisor equality and closedness of $S$
give the required containment.  Suppose $x$ belongs to the zero locus
and to $E$, but not to $S$.  In a neighbourhood missing $S$, choose an
actual nondegenerate local generator $f$ for $R$ and genuine flattening
coordinates for the smooth hypersurface $E$.  Nondegeneracy makes the
coordinate pullback of $f$ a nonzero germ.  The actual member outside
$I_E$ excludes the exceptional coordinate factor, since the complete
principal stalk is preserved in both directions by these holomorphic
coordinates.  At any nearby point outside $E$, the actual ideal is the
unit ideal, because that point is outside $S$.  Hence the zero germ of
$f$ is contained in the exceptional coordinate hyperplane.  The preceding
nonexceptional equation lemma makes $f(x)\ne0$, contrary to $x$ being
in the original zero locus.  This proves the converse containment.
\end{proof}

\begin{theorem}[The entire reduced ideal across a smooth divisor]
\label{hir:invertible-reduced-ideal-across-smooth-divisor}
\uses{hir:invertible-ideal-support-across-smooth-divisor,hir:reduced-germ-across-coordinate-divisor,hir:squarefree-equation-full-vanishing-ideal}
\lean{Hironaka.HironakaHypersurfaceStrictIdeal.eq_vanishingIdeal_of_reduced_off_smooth_divisor}
\leanok
Under the preceding geometric hypotheses, $R=I_S$ as entire holomorphic
ideal sheaves, including all stalks on $E$.
\end{theorem}
\begin{proof}
The support equality and coherence make every actual member of $R$ vanish
on $S$ near its base point, giving $R\subseteq I_S$.  Away from $E$ the
other inclusion is already the actual off-divisor stalk equality.
At $x\in E$, choose the same genuine generator and flattening coordinates.
The original whole-ideal witness excludes the exceptional coordinate from
the pulled-back generator $f$.  At every nearby point outside that
coordinate hyperplane, the principal stalk of $f$ is the actual vanishing
ideal of the original $S$ in these coordinates.  It is radical: a positive
power vanishing on $S$ forces the germ itself to vanish there.
Both the generator and arbitrary test germs transport through the same
actual local holomorphic inverse; thus this is the whole principal stalk,
not merely its zero locus.  The coordinate-divisor reducedness theorem
therefore proves that $f$ is squarefree at the original point $x$.
Its entire zero-germ vanishing ideal is $(f)$.  By the support equality,
an arbitrary germ in $I_{S,x}$ vanishes on that same zero germ.  It is
therefore a holomorphic multiple of $f$ in the coordinates.  Pull the
multiple back through the actual coordinate inverse.  This gives membership
in the original complete stalk $R_x$.  Since this holds at every $x$, the
two ideal sheaves are equal.
\end{proof}

\begin{lemma}[The off-centre dictionary for the actual strict transform]
\label{hir:controlled-hypersurface-off-centre-strict-stalk}
\uses{hir:controlled-hypersurface-off-centre-stalk}
\lean{Hironaka.HironakaHypersurfaceStrictIdeal.controlledTransform_vanishingIdeal_strict_mem_iff_of_notMem}
\leanok
For a closed $W$ and an actual blowup $\pi$ along $C$, at each
$y\notin\pi^{-1}C$ the whole controlled ideal stalk, with any mark,
equals $I_{\operatorname{strictTransform}(\pi,W),y}$.
\end{lemma}
\begin{proof}
The exceptional locus lies in $\pi^{-1}C$.  On the open complement of
that preimage the closed strict transform and $\pi^{-1}W$ have the same
membership germ: inverse-image points already belong to the set whose
closure defines the strict transform, and the closed strict transform
lies in $\pi^{-1}W$.  Their whole vanishing ideal stalks agree.  Apply
the actual off-centre controlled-ideal dictionary above.
\end{proof}

\begin{theorem}[Actual strict ideal after a maximum-order hypersurface blowup]
\label{hir:actual-max-order-hypersurface-strict-ideal}
\uses{hir:invertible-controlled-hypersurface-ideal,hir:controlled-hypersurface-off-centre-strict-stalk,hir:invertible-reduced-ideal-across-smooth-divisor,hir:actual-fresh-maximal-multiplicity-centre}
\lean{Hironaka.HironakaHypersurfaceStrictIdeal.controlledTransform_vanishingIdeal_eq_strictTransform}
\leanok
Let $W\subset X$ be an analytic hypersurface in an $(m+1)$-dimensional
complex manifold, with $W\subset\overline{\operatorname{Reg}_mW}$ and
actual invertible vanishing ideal.  Let $E$ be an SNC boundary and
$\mu\ge2$ an upper bound on the actual orders of $I_W$ on $W$.
For any actual admissible centre of $(X,X,E,I_W,\mu)$ and its actual
blowup $\pi:Y\to X$, the controlled ideal is the entire reduced
vanishing ideal of $S=\operatorname{strictTransform}(\pi,W)$.
\end{theorem}
\begin{proof}
Admissibility puts the centre in the order-$\mu$ locus.  Order is one
at each regular point, so the centre lies in $W\setminus\operatorname{Reg}_mW$.
The original upper bound and this lower bound make the actual order
exactly $\mu$ at every centre point.  Hadamard centre divisibility
gives $I_W\subset I_C^\mu$, and the invertibility theorem gives an
actual invertible controlled ideal $R$.

The updated boundary, including $D=\pi^{-1}C$, is SNC, hence its
original last member $D$ is a smooth hypersurface.  Since the regular
locus of $W$ avoids $C$ and is dense, $W\subset\overline{W\setminus C}$.
The geometric strict transform $\overline{\pi^{-1}(W\setminus C)}$
is therefore the actual library strict transform $S$.  Its defining
dense subset lies in $S\setminus D$, so $S\subset\overline{S\setminus D}$.
The actual off-centre dictionary supplies $R_y=I_{S,y}$ off $D$.

Exact order $\mu$ excludes order $\mu+1$ at the image of every point
of $D$.  The actual exceptional pushforward theorem consequently
produces, at each such point, an element of the controlled ideal which
is not in $I_D$.  This is a whole germ membership assertion, including
at points where its value vanishes.  The reducedness-across-the-divisor
theorem now applies with all its inputs internally supplied, and gives
$R=I_S$ as complete ideal sheaves.
\end{proof}

\begin{theorem}[Internally produced next reduced hypersurface ideal]
\label{hir:fresh-max-order-blowup-strict-ideal}
\uses{hir:actual-max-order-hypersurface-strict-ideal,hir:actual-fresh-maximal-multiplicity-centre}
\lean{Hironaka.HironakaHypersurfaceStrictIdeal.exists_fresh_maximal_order_blowup_strictIdeal}
\leanok
For a compact analytic hypersurface $W$ satisfying the preceding density
and invertibility conditions, an SNC boundary, and an actual point of
order at least two, there exist an internally chosen finite maximum
order $\mu\ge2$, a nonempty admissible centre in the current singular
locus with constant order $\mu$, and an actual blowup such that the
next current hypersurface is analytic and compact, has invertible
vanishing ideal, and satisfies
\[
 \pi^*I_W=I_{\pi^{-1}C}^{\mu}I_S,\qquad
 \operatorname{controlledTransform}(I_W,\mu)=I_S.
\]
The transformed old boundary together with the actual new divisor is SNC.
\end{theorem}
\begin{proof}
The actual compact upper semicontinuous reduced order function attains
its finite maximum.  Use the actual canonical maximum locus of the
marked invariant to produce its nonempty admissible centre; the
maximum-order centre construction supplies the global bound and exact
order there.  Resolve the actual smooth centre by the actual blowup
constructor.  Apply the preceding whole-ideal theorem; its
invertibility and the ordinary coherent controlled-product identity
give the displayed factorization with the actual $I_S$.
Analyticity and compactness of $S$ are the actual blowup strict-transform
theorems.  The actual admissible transform supplies the updated SNC list.
\end{proof}

The first lemma's radicality is local.  In the actual blowup producers
it follows from the original reduced ideal through the blowup's
off-centre biholomorphism; no reduced-transform certificate is assumed.
The arbitrary pure codimension-one local equation producer and the finite
restarted multiplicity descent remain separate open requirements.
The general G8 assertion in higher codimension remains open.
